% TODO[HLA-FIX 2026-06-27]: PredIG 全局显著性已失效(max-pool rho 0.198->0.104 p=0.343 ns; mean-agg 0.280->0.188 p=0.084 ns)，剔污染患者 P101/P102 后不存活；IMPROVE 仍稳健显著(0.226 p=0.037)，TSCAPE 翻为显著负(-0.230 p=0.033)，deepHLApan 另有 merge bug 双重不可信；per-patient 头条数字(PRIME/deepHLApan)待 Phase B 正确等位重推理后更新。正文数字暂不动(投稿=拍板点 + P102 等位待袁老师确认)。详见 04_LOG Entry HLA-FIX。
\section{Discussion}
\label{sec:discussion}

\subsection{Existing tools do not quantify response magnitude}
\label{sec:disc-gap}

The benchmark gives a direct, empirical answer to the question that motivated it:
when the scores of current immunogenicity predictors are repurposed as continuous
estimators of T-cell response magnitude, they track it poorly. Across the eight
tools, every absolute Spearman correlation with ELISpot magnitude falls below
$0.33$, and only IMPROVE and PredIG reach significance, and then only under
% XXX_HLA-FIX[投稿前必改]: 修复后此句不再成立——剔 P101/P102 后 PredIG 全局显著性失效(p=0.343 ns)，"only IMPROVE and PredIG" 须改为 "only IMPROVE"。数字暂不动=拍板点，见文件头 TODO + 04_LOG Entry HLA-FIX。
particular aggregation choices (Section~\ref{sec:results}). This is consistent
with, and quantifies, the gap we identified in the literature survey: no surveyed
method is designed to regress magnitude, and several discard the graded signal
their own training data contain. IMPROVE and PredIG are trained on data that
distinguish graded response levels but are binarized before use
\cite{IMPROVE2024,PredIG2025}, and other tools collapse graded labels into a
single recognition class. The benchmark therefore turns a documentary
observation---that the field reports binary calls---into a measured one: even the
best-performing current scores explain only a small fraction of the variance in
how strongly patients respond. We read the convergence of the survey evidence and
the empirical correlations as showing that the magnitude gap is real rather than
an artifact of how individual papers report their results.

A weak correlation, however, is not automatically a deficiency to be engineered
away. The same result is compatible with two readings---tools that could do
better but were never asked to, or a target that is intrinsically hard to predict
from the available inputs---and the next subsection argues that both readings hold
simultaneously.

\subsection{A biological ceiling bounds what sequence-level prediction can achieve}
\label{sec:disc-ceiling}

There is a structural reason to expect that magnitude prediction from peptide and
allele alone is bounded well short of perfect. The magnitude of a polyclonal
T-cell response is, to a first approximation, the product of a recruitment-and-
expansion chain whose leading term is the number of naive precursor T cells
specific for the epitope; in mouse models, naive precursor frequency is reported
to be a major determinant of the magnitude of the elicited
response~\cite{Jenkins2012,Moon2007}. The quantitative form of this dependence in
human neoantigen-specific CD8 responses has not been established and the
extrapolation across species and T-cell subset is uncertain, but the qualitative
direction---that a host-specific, sequence-invisible quantity gates response
size---is expected to carry over. That quantity is a property of the
host's T-cell repertoire and thymic selection history, varies substantially
across donors for the same peptide and allele, and is in principle unobservable
from a peptide-and-HLA input. Writing the explainable variance as a ratio,
\begin{equation}
  \rho_{\max}^2 \;\le\; \frac{\operatorname{Var}\!\big(\mathbb{E}[Y \mid X]\big)}
                              {\operatorname{Var}(Y)},
  \label{eq:ceiling}
\end{equation}
where $Y$ is the measured magnitude and $X$ the peptide-and-allele input, the
donor-specific precursor and repertoire terms sit in the numerator's complement
and cannot be recovered no matter how much sequence data are supplied. A
first-order analysis along these lines places the attainable rank correlation at
roughly $\rho_{\max}\approx 0.4$ to $0.6$; we present this figure as a positional
estimate rather than a measured bound, since it rests on biological variance
parameters we did not measure directly. Read against this ceiling, the best
magnitude correlation under the locked primary protocol (IMPROVE,
$\rho = 0.2434$) already reaches roughly half of the lower end of the plausible
maximum; even the per-tool upper envelope obtained by allowing each tool its own
best-of-nine operating point---a descriptive figure we do not treat as a fair
ranking (Section~\ref{sec:fairness})---tops out near $\rho\approx 0.32$ for IMPROVE,
on the order of half to two-thirds of the ceiling. Either way the headroom for
pure sequence-level magnitude regression appears modest rather than large.

We take a deliberate position on what this implies for tool design. If the
ceiling on pure peptide-and-HLA inputs is genuinely in the middle of the range,
then the value of a continuous-magnitude model should not be advertised as
high-precision regression. It should instead be framed as recovering
ranking-relevant signal for clinical top-$K$ prioritization---where only a handful
of peptides can be synthesized into a vaccine and the quality of the ordering is
what matters---which binary tools discard by construction. Breaking past the
ceiling would require feeding donor-specific information (HLA typing, T-cell
receptor repertoire, or precursor estimates) rather than a more elaborate
sequence model.

\subsection{Per-patient evaluation should be the default}
\label{sec:disc-perpatient}

A methodological lesson generalizes beyond the tools we ranked. Pooling all
peptides across patients into a single global correlation can hide a tool's
ability to order candidates \emph{within} an individual patient, which is the
quantity that actually governs vaccine selection. Several tools whose pooled
global correlation is indistinguishable from noise show within-patient ranking
ability in a subset of patients---alongside substantial dispersion---once the
correlation is computed per patient and then summarized: PRIME moves from a global
Spearman of $0.116$ to a per-patient median coefficient of $0.386$ (and a
Fisher-$z$ weighted mean of $0.253$), and because its best-binder scores are
essentially uncorrelated with peptide length, this gain reflects a genuine
within-patient reordering rather than a length artifact
(Section~\ref{sec:results}). A per-patient jump must nonetheless be checked against
aggregation-level confounds before it is read as capability: deepHLApan shows the
numerically largest shift (global $0.042$ to a per-patient median of $0.402$), but
under best-binder aggregation its pooled scores correlate with peptide length
($\rho \approx 0.57$) and the signal collapses to $\approx 0$ under
length-insensitive pooling, so we treat its jump as a cautionary case rather than
as evidence of within-patient ranking ability. Because of the wide individual
heterogeneity ($\rho_i$ spanning roughly $-0.43$ to $+0.81$ for the most dispersed
tools) we read the per-patient aggregates as descriptive summaries of central
tendency rather than as estimates of a single common correlation, and we anchor the
qualitative claim on the median together with the reported spread. The global metric is not wrong, but it answers a
different question---can a single score threshold separate responders across an
entire cohort---than the clinically relevant one. We therefore recommend that
benchmarks of magnitude or ranking ability compute the metric per patient before
aggregating, mirroring the cohort-level paradigm of the clinical neoantigen
literature \cite{TESLA2020,Muller2023}, and report the across-patient spread
rather than a single pooled number.

\subsection{Limitations}
\label{sec:disc-limits}

Several limitations bound our conclusions and we state them plainly. First, the
primary cohort is small and imbalanced: DS2 contains $101$ peptides with only
$11$ negatives at the SFC $>0$ threshold and $K=9$ patients with $5$ to $16$
peptides each, so every per-patient aggregate carries a wide confidence interval
and the per-patient point estimates should not be over-read. We report bootstrap
and Fisher-$z$ intervals throughout for this reason, and we draw conclusions from
the consistency of the pattern across aggregation choices rather than from any
single coefficient. Second, our results rest on a single benchmark cohort;
external replication on independent ELISpot cohorts is needed before the ranking
behavior can be generalized. Third, several tools were trained on IEDB-derived
data that may overlap with the benchmark peptides, which would bias those tools
optimistically; we note that this works \emph{against} our central claim rather
than for it, since even with any such advantage no tool achieves a strong
magnitude correlation. Fourth, three tools are evaluated in non-canonical
configurations that we carry through every table: NeoTImmuML is a model retrained
by us because official weights are unavailable (denoted NeoTImmuML$^\star$),
pTuneos is run through its Pre\&RecNeo sub-model, and IMPROVE is run with its
expression feature degraded; our statements about these tools refer to the
configuration actually evaluated, not necessarily to the original published
models. Because a degraded or retrained configuration can only lower a tool's
apparent performance relative to its full published form, our central claim---that
current tools do not recover response magnitude---is, for these three tools, a
conservative statement rather than an upper bound on what the original models could
achieve; it remains directly supported by the five tools evaluated in their
intended configurations. Finally, the expanding second wave of predictors and the tools governed by
restrictive academic data-use terms are integrated into the pipeline but their
results are withheld from the main conclusions, which do not depend on them.

\subsection{Future work}
\label{sec:disc-future}

The natural sequel to this benchmark is a predictor built to occupy the L4 rung
directly: a model that regresses continuous response magnitude against ELISpot
SFC rather than emitting a binary call. The benchmark supports such an effort in
two concrete ways. It establishes that the gap is real and currently unfilled,
which is the empirical justification for attempting it, and it provides a
reproducible evaluation protocol---harmonized inputs, the best-binder aggregation,
the per-patient analysis, and the bootstrap intervals---against which any
candidate magnitude predictor can be measured on equal terms. Consistent with the
ceiling argued in Section~\ref{sec:disc-ceiling}, we expect the realistic payoff
of such a model to lie in improved clinical top-$K$ ranking rather than in
high-accuracy regression, and we regard the incorporation of donor-specific
information as the most promising, if speculative, route to pushing past the
sequence-level ceiling.
